% Numbers are macros from tables/numbers.tex. Statements that depend on a rule outcome name the rule.

\paragraph{The estimator, not the geography.} What separates the explainers in our results is where each one queries
the model. LIME's quartile sampler produces rows that no real home resembles; a detector separates them perfectly,
and LIME then explains the innocuous model while the deployed model keeps penalising almost every real home. The
Shapley explainers are harder to fool for two reasons that have nothing to do with space. Efficiency ties their
attributions to the served prediction at the explained home, so the penalty cannot leave the explanation while that
home is served by the biased model; and their coalition hybrids include, at both ends of the size range, points that
differ from a real home in one or two values and that a detector cannot reject without rejecting real homes.
GeoShapley moves the two coordinates together, which keeps every hybrid at a real location and makes its queries
slightly harder to detect, but its other players are still recombined across homes, and the middle coalitions are as
far from the data as KernelSHAP's. Under rule C2 we find no detectable difference between the two in what the
attacker can hide. Treating location carefully is valuable for interpretation; it is not, on its own, a defence. The
question to ask of any explainer used in an audit is the one the routing curves answer: what share of its queries
looks like real data?

\paragraph{What partial hiding means for an audit.} ``Partly fooled'' is not ``safe''. Against the matched scaffold
KernelSHAP and GeoShapley keep between half and three quarters of the penalty, and an auditor who knows what to look
for would still find the trigger among the leading features for about half of the penalised homes. But an audit is
rarely a search for one named variable, and an attacker who accepts lower fidelity can do much better: at 50\%
fidelity the scaffold still charges half of the targeted homes the full penalty while KernelSHAP shows
\FKernelSHAPMatFiftyKept\ of it. The frontier in Figure~\ref{fig:F3_frontier} turns this into a quantity a regulator
can reason about: for a given explainer, how much bias can be hidden at what cost in bias actually delivered.

\paragraph{GeoShapley under attack.} For geographers the more specific lesson is about how GeoShapley fails. Under
the scaffold the intrinsic location effect on penalised homes shrinks towards zero and the feature-by-location
interactions roughly double (rule C3). A GEO map read on its own would understate the penalty, and the interaction
terms would invite a substantive story --- features that matter differently in different places --- that is in fact
an artefact of the attacker's routing. We did not set out to build a detector and do not evaluate one, but inflated
interactions without a corresponding change in the model's fit are a pattern worth examining in future work.

\paragraph{Reconciling with earlier results.} \citet{slack2020fooling} found LIME easier to fool than SHAP on
tabular benchmarks; we find the same ordering in a housing market. An earlier version of this project reported the
opposite, that LIME resists the attack in a geographic setting. That result came from the generic detector, which is
trained on Gaussian noise and lets \FLIMEGenLeak\% of LIME's real queries through. Replications of scaffold attacks
should match the detector to the explainer's actual sampler, or report how many queries reach the biased model.

\paragraph{Changing the audit settings.} \todo{after the robustness runs: which audit-setting changes, if any, move
the share kept (Table~\ref{tab:robust}); whether an auditor gains anything from a larger or $k$-means background or
from LIME's Gaussian sampler; keep the claim to what rule R and the intervals support.}

\paragraph{Looking at outcomes.} The penalty that the scaffold hides from explainers is plainly visible in the
residuals of the model's own predictions: their local mean separates penalised homes with AUC
\EsevenAucLagScaffoldTen\ even at a 21\% penalty, using tools --- local means, Moran's $I$, LISA --- that are standard
in geography. This check needs observed outcomes, loses power for penalties of a few percent, and we have not tested
it against an attacker who anticipates it. It complements explanation-based audits rather than replacing them, and it
reflects a broader point: for spatial bias, an audit that looks at what a model does to places can be more robust than
one that asks the model to explain itself.
